\documentclass[10pt,a4paper]{amsart}
\usepackage[margin=19mm]{geometry}
\usepackage{amsmath,amssymb,mathtools}
\usepackage[scr=rsfs,cal=euler]{mathalfa}
\usepackage{xfrac}
\usepackage[unicode,hidelinks,bookmarks=false]{hyperref}
\newcommand{\ie}{i.e.\ }
\newcommand{\cf}{cf.\ }
\newcommand{\resp}{respectively\ }
\newcommand{\Subsection}{\subsection}
\providecommand{\nobreakdash}{\nobreak\hbox{-}}
\setlength{\emergencystretch}{2em}
\multlinegap0pt
\usepackage{bm}
\usepackage{bbm}
\usepackage{dsfont}
\usepackage{xspace}
\usepackage{cancel}
\usepackage{mathtools}
\usepackage{amsthm}
\makeatletter
\@ifundefined{theorem}{\newtheorem{theorem}{Theorem}}{}
\@ifundefined{lemma}{\newtheorem{lemma}[theorem]{Lemma}}{}
\makeatother
\newcommand{\F}{\mathbb{F}}
\title[$2$-quasi-perfect Lee codes and abelian Ramanujan graphs: a ]{$2$-quasi-perfect Lee codes and abelian Ramanujan graphs: a new construction and relationship}
\author{Shohei Satake}
\date{Source: arXiv:2601.12393v3 (CC BY 4.0)}
\begin{document}
\maketitle

\noindent\emph{Source and licence.} $2$-quasi-perfect Lee codes and abelian Ramanujan graphs: a new construction and relationship. Version arXiv:2601.12393v3 (CC BY 4.0), distributed under the Creative Commons Attribution 4.0 International licence, as recorded in the arXiv licence metadata for this exact version and shown on the abstract page https://arxiv.org/abs/2601.12393v3. Full rights notice: licenses/arxiv-2601.12393-CC-BY-4.0.txt.\par\smallskip

\noindent\textit{Editorial scope.} Counted unit: the Lemma whose source label is \texttt{lem-count} in the article (source0.tex lines 276--291, source label \texttt{lem-count}), delivered together with the article's own complete original proof of it (source0.tex lines 293--327, one \texttt{proof} environment of 35 lines). No mathematics is written, repaired or re-derived: statement and proof are the article's own text line for line. Required context retained uncounted: none. Every definition, display and label of those lines is retained unchanged. Omitted from this package and not redistributed: 1--275, 292--292, 328--770. Nothing inside a retained excerpt is omitted or altered: every retained body is delivered exactly as the article's lines are written, so no omission receipt of any kind accompanies this package. Citations: the retained excerpts carry no citation command at all, so no bibliography entry of the article is transcribed and no reference is redistributed. Reference adaptation: none was needed and none was made; every reference inside the delivered unit resolves inside it, so no unresolved reference or citation is produced. Printed slips kept literal: no printed word, symbol, hypothesis or number of the counted statement or of its proof was altered. Layout and class: the article's own class is \texttt{amsart} with packages including geometry, graphicx, amsmath, amsthm, amssymb, amsfonts, mathtools, hyperref; the wrapper is the lane's pinned \texttt{amsart} editorial wrapper at 10pt on A4, which restates the theorem-like environments the retained text needs and adds the article's own macro definitions that the retained text needs (\texttt{\textbackslash F}), with extra wrapper packages (bm, bbm, dsfont, xspace, cancel, mathtools). The delivered numbering is the unit's own. Assets: no retained span contains a figure environment, a graphics-inclusion command, a PDF-page inclusion, a table or a TikZ picture; no image file is copied into this package, the package declares no asset dependency and no image file is redistributed with it. Licence: version arXiv:2601.12393v3 of this article is distributed under the Creative Commons Attribution 4.0 International licence, as recorded in the arXiv licence metadata for this exact version and shown on the abstract page https://arxiv.org/abs/2601.12393v3. A full rights notice is delivered at licenses/arxiv-2601.12393-CC-BY-4.0.txt. Characters: every retained excerpt is pure ASCII, so the delivered engine, XeLaTeX with its default Latin Modern text font, reports no missing character.
\begin{lemma}
\label{lem-count}
For any prime $p \geq 5$ and power $q=p^k$ with $k\geq 1$, 
we have 
\begin{align*}
    \#\bigcup_{i=0}^{2} H_q^{(i)}=2\cdot \Bigl(\frac{q-1}{2} \Bigr)^2+2\cdot \Big(\frac{q-1}{2} \Big)+ 1
\end{align*}
% \begin{enumerate}
%     \item[$(1)$] $\#H_q^{(2)}=\frac{(q-1)^2}{2}+1$,
    
%     \item[$(2)$] $\{0\}=\frac{(q-1)^2}{2}+1$,
% \end{enumerate}
% \begin{align*}
%     \#H_q^{(2)}=\frac{(q-1)^2}{2}+1
% \end{align*}
\end{lemma}

\begin{proof}
First we show that $\#H_q^{(2)}=\frac{(q-1)^2}{2}+1$.
It suffices to count $(a, b)\in \F_q^2$ such that there are $x, y\in \F_q^*$ with
\begin{align}
\label{eq-count}
  \begin{cases}
    a=x+y  \\
    b=x^3+y^3
  \end{cases}
\end{align}
Note that $a=0$ if and only if $y=-x$, equivalent to $b=0$, implying that there are exactly $\frac{q-1}{2}$ unordered pair $\{x, y\}$ satisfying \eqref{eq-count} with $(a, b)=(0, 0)$. 
If $a\neq 0$, then $(a, b)$ corresponds to an unordered pair $\{x, y\}$ satisfying \eqref{eq-count} and $x+y\neq 0$ by a one-to-one mapping. 
Indeed, by \eqref{eq-count}, we have $xy=\frac{a^3-b}{3a}$, and hence $x, y$ must be roots of the quadratic equation $t^2-at+\frac{a^3-b}{3a}=0$. 
The total number of unordered pairs $\{x,y\}$ from $\mathbb{F}_q^*$ (allowing $x=y$) is equivalent to choosing 2 elements from $q-1$ elements with replacement, which is $\binom{(q-1)+2-1}{2} = \binom{q}{2} = \frac{q(q-1)}{2}$.
%The pairs satisfying $x+y=0$ are exactly $\{x, -x\}$ for $x \in \mathbb{F}_q^*$, and there are exactly $\frac{q-1}{2}$ such pairs.
Thus, the number of unordered pairs with $x+y \neq 0$ is:
%Now we obtain $\#H_q^{(2)}$ as the number of $\{x, y\}$ with $x+y\neq 0$ is
\begin{align*}
    \frac{q(q-1)}{2} - \frac{q-1}{2} = \frac{(q-1)^2}{2}.
\end{align*}
% \begin{align*}
%     {}_{q-1}H_2-\frac{q-1}{2}=\binom{(q-1)+2-1}{2}-\frac{q-1}{2}=\frac{(q-1)^2}{2}.
% \end{align*}
Hence, we conclude that $\#H_q^{(2)} = \frac{(q-1)^2}{2}+1$ by adding $(0,0)$.
%Notice that $(0, 0) \in H_q^{(2)}$, implying that $H_q^{(0)} \subset H_q^{(2)}$.
The proof completes after showing that $H_q \cap H_q^{(2)}=\emptyset$ as $H_q^{(0)} \subset H_q^{(2)}$ and $\#H_q^{(1)}=\#H_q=q-1$.
Suppose that there exists $(a, a^3) \in H_q$ such that there are $x, y \in \F_q^*$ with 
\begin{align*}
  \begin{cases}
    a=x+y  \\
    a^3=x^3+y^3
  \end{cases}
\end{align*}
This implies that $3axy=0$, which is impossible as $x, y, a \neq 0$ and $p\geq 5$.
\end{proof}

\end{document}
